% !PRINTED 269
\weberpage{269}
%<<CONT
sign changes, whereas for negative discriminant, where $a$ is also negative, \gesperrt{one} sign change takes place.

So long as $a$ lies in the interval (9), our equation therefore has two imaginary roots; if it lies outside this interval, four imaginary roots, never four real roots.

A purely numerical example is offered by the equation, drawn from the complex multiplication of the elliptic functions,
\[
x^5 - x^3 - 2x^2 - 2x - 1 = 0.
\]
For the determinant of the Bezoutiant we obtain from \S.\,73
\[
\begin{vmatrix}
\dfrac{-4}{5}, & \dfrac{-3}{5}, & \dfrac{4}{5}, & 5 \\[6pt]
\dfrac{-3}{5}, & \dfrac{4}{5}, & \dfrac{33}{5}, & 8 \\[6pt]
\dfrac{4}{5}, & \dfrac{33}{5}, & \dfrac{46}{5}, & 6 \\[6pt]
5, & 8, & 6, & 2
\end{vmatrix}
\]
the chain of principal subdeterminants
\[
\frac{47^2}{5}, \; \frac{94}{5}, \; -1, \; \frac{-4}{5}, \; 1,
\]
so that the discriminant of our equation is $47^2$. The equation therefore has two pairs of imaginary roots and one real root. The second term of this chain, $94 : 5$, need not be computed if one knows the discriminant; indeed, the mere knowledge of the negative sign in the penultimate term already suffices to reach the foregoing decision on the reality of the roots.
